% ============================================================
% sections/s03_task_allocation.tex —— 3 滚动时域选择性任务分配方法
%
% 【术语红线（写作要点，注释不输出；执笔本节前先读）】
% 1) 本节「阶段一：选择性任务接纳 / 阶段二：多机巢协同任务分配」是求解
%    框架的划分，与 §2.3 的两层目标 (P1)/(P2) 不是同一划分，禁止互指，
%    更不得写成「先解 (P1) 求 Y_t^*、再解 (P2)」。实现只用一个加权标量 f
%    （覆盖门禁 + 支配权重 W_skip/W_hard 把覆盖优先压进高层罚项），
%    统一表述为「以支配权重实现覆盖优先」。
% 2) 代码里三个 lexicographic* 同名不同物，均不得当作论文「字典序」：
%    - config.py 旧名 lexicographic_coverage（现 coverage_audit）：只是
%      W_skip 罚层 / 审计分账开关，config.py:127 明写「不是字典序口径」；
%    - config.py 的 lexicographic 死开关：从不改变行为，已失效移除；
%    - _select_by_lexicographic_marginal（solvers/participating.py）：
%      阶段一选机的边际比较 (Δ可服务数, −Δ覆盖代价)，是局部启发式判据，
%      不是全周期目标层次。
%    论文「字典序」= §2.3 数学规划的目标优先序（先交付数、后成本）。
% 3) 符号表 σ_t 写「估计口径见第 3 节」，本节须落 cap(t)/n_sel 公式
%    （methods/selection/capability.py；文件头已标「论文 3.3」）：
%    cap(t)=max(min_cap, ⌊B(t)/c̃⌋)，n_sel=min(⌈max_ratio·K⌉, cap(t))；
%    σ_t 是接纳数上界而非等式，估计偏高时自动退化到真实可达上限。
%    在本节补写前，符号表那句「见第 3 节」仍是悬空引用。
% ============================================================

\section{滚动时域选择性任务分配方法}

\subsection{总体框架}

\placeholder{}

\subsection{阶段一：选择性任务接纳机制}

\placeholder{}

\subsection{阶段二：多机巢协同任务分配}


\subsection{滚动闭环与状态继承}

\placeholder{}
